%%=============================================================================
%% Methodologie
%%=============================================================================

\chapter{\IfLanguageName{dutch}{Methodologie}{Methodology}}%
\label{ch:methodologie}

Om het onderzoek zo succesvol mogelijk te laten verlopen, werd ervoor gekozen het op te splitsen in zes fasen. De eerste fase is een probleemverkenning, waarin een literatuurstudie wordt uitgevoerd om meer inzicht te krijgen in het probleemdomein. Vervolgens wordt een requirements-analyse uitgevoerd, resulterend in een lijst van beoordelingscriteria. Na de requirements-analyse volgt een inventarisatie van het huidige digitale (of analoge) aanbod van burn-outbestrijdingsproduct- en, de zogenaamde longlist-fase. Eenmaal deze producten in kaart zijn gebracht, worden ze teruggebracht tot een shortlist en beoordeeld aan de hand van de eerder opgestelde beoordelingscriteria. Vervolgens is het tijd om de \acrlong{poc} op te stellen. In deze fase worden mock-ups ontwikkeld, de benodigde tools geselecteerd en uiteindelijk een werkend prototype gerealiseerd. Tot slot wordt in de conclusiefase het succes van dit werkende prototype bepaald door middel van een heuristische evaluatie.

\section{Probleemverkenning}
\label{sec:probleemverkenning}

\vspace{1em}

Gedurende de eerste fase is het de bedoeling een diepgaand begrip te krijgen van het probleemdomein aan de hand van een \textbf{literatuurstudie}. Voor dit onderzoek is het essentieel te begrijpen wat een burn-out precies inhoudt, hoe dit speelt bij de huidige generatie twintigers, welke rol digitalisering en een verhoogde cognitieve belasting hierbij kunnen spelen, en welke preventie- en behandelingsstrategieën momenteel beschikbaar zijn. 

\vspace{1em}

Vervolgens wordt het huidige aanbod aan digitale oplossingen voor mentale gezondheid onderzocht, met aandacht voor hun effectiviteit en mogelijke tekortkomingen. Daarbij worden ook UX- en ontwerpprincipes die bijdragen aan mentale rust onder de loep genomen. Tot slot wordt gekeken naar de bestaande technologieën in de app-ontwikkelingswereld, zodat de meest geschikte technologie kan worden geselecteerd voor de ontwikkeling van de proof-of-concept.

\vspace{1em}

Het resultaat van deze fase is een overzicht van de oorzaken van burn-out bij twintigers, hun belangrijkste risicofactoren, en een eerste inventaris van bestaande \newline preventie- en behandelingsstrategieën.

\section{Requirements-analyse}
\label{sec:req-analyse}

\vspace{1em}

Tijdens de requirements-analyse wordt bepaald wat de exacte noden van de doelgroep zijn door middel van het opstellen van functionele en niet-functionele requirements. Deze kunnen enerzijds worden afgeleid uit de reeds verrichte literatuurstudie en anderzijds worden aangevuld door één of meerdere interviews met psychologen. De focus binnen deze studie ligt op het in kaart brengen van effectieve behandelingsmethoden en gebruikerseisen, welke op hun beurt als requirements kunnen worden gedefinieerd.

\vspace{1em}

De requirements worden opgesteld volgens de MoSCoW-methode (Must, Should, Could, Won’t). Deze aanpak maakt het mogelijk om prioriteiten helder te definiëren en de kernfunctionaliteiten concreet te maken. Aan het einde van deze fase ligt er een geprioriteerde lijst van functionele en niet-functionele requirements klaar, die als basis dient voor de \acrlong{poc}.

\section{Longlist}
\label{sec:long-list}

\vspace{1em}

Na de requirements-analyse volgt de inventarisatie van het huidige digitale en analoge aanbod van burn-outbestrijdingsproducten, de longlist-fase. In deze fase worden bestaande alternatieven, zoals apps, websites, wearables en andere concepten, verzameld en geanalyseerd. Alleen oplossingen met een bruikbaar concept dat eventueel digitaal kan worden toegepast, worden opgenomen. De alternatieven worden opgespoord via online bronnen, app-stores en bestaande studies. Elk alternatief wordt vervolgens een eerste keer getoetst aan de opgestelde requirements, zonder dat er al een definitieve beoordeling plaatsvindt. Het resultaat is een overzicht van alle mogelijke alternatieven, inclusief bronvermelding, dat als basis dient voor de volgende fase: de shortlist.

\section{Shortlist}
\label{sec:short-list}

\vspace{1em}

In de shortlist-fase wordt uit de longlist een selectie gemaakt van de meest veelbelovende opties die verder onderzocht zullen worden. Een overzichtelijke matrix ondersteunt de beoordeling van de alternatieven ten opzichte van de gestelde requirements, zodat duidelijk wordt welke oplossingen het beste aansluiten bij de verzamelde criteria en de meest bruikbare componenten bevatten. Aan het einde van deze fase wordt een keuze gemaakt voor één of enkele alternatieven, inclusief de bijbehorende componenten, die meegenomen zullen worden in de ontwikkeling van de \acrlong{poc}.

\section{Proof-of-Concept}
\label{sec:poc}

\vspace{1em}

Tijdens de \acrlong{poc}-fase ligt de focus op het ontwikkelen van een werkend prototype. De ontwikkeling start met het ontwerpen van mock-ups, met speciale aandacht voor mentale rust en minimale cognitieve belasting. Tegelijkertijd wordt bepaald welke app-ontwikkeltechnologie het meest geschikt is. Zodra ontwerp en technologie zijn vastgesteld, volgt de daadwerkelijke uitwerking van het prototype. Dit prototype wordt regelmatig afgetoetst bij de copromotor of andere psychologen om validatie vanuit een professioneel standpunt te waarborgen. Eindgebruikers, met een burn-out of het risico daarop, worden niet betrokken bij deze evaluatie, aangezien zij geen objectief oordeel kunnen geven over de effectiviteit van de oplossing op korte termijn.

\section{Conclusie}
\label{sec:conclusie}

\vspace{1em}

Tot slot, in de conclusiefase, wordt het succes van het werkende prototype beoordeeld aan de hand van een heuristische evaluatie, gebaseerd op de eerder opgestelde evaluatiecriteria en requirements. Deze evaluatie wordt zowel zelfstandig uitgevoerd als door experts, zoals de copromotor, om een professioneel perspectief te waarborgen. Eventuele hiaten of tekortkomingen worden geïdentificeerd, en er worden waar nodig aanbevelingen gedaan voor toekomstig onderzoek. 

\vspace{1em}

Het concrete resultaat van deze fase is een afgeronde analyse van de effectiviteit van de gekozen oplossing, inclusief een overzicht van relevante requirements die ook van waarde kunnen zijn voor andere burn-outgerelateerde IT-projecten. Op deze manier wordt zowel de praktische toepasbaarheid als de wetenschappelijke onderbouwing van het prototype duidelijk gemaakt.

%% TODO: In dit hoofstuk geef je een korte toelichting over hoe je te werk bent
%% gegaan. Verdeel je onderzoek in grote fasen, en licht in elke fase toe wat
%% de doelstelling was, welke deliverables daar uit gekomen zijn, en welke
%% onderzoeksmethoden je daarbij toegepast hebt. Verantwoord waarom je
%% op deze manier te werk gegaan bent.
%% 
%% Voorbeelden van zulke fasen zijn: literatuurstudie, opstellen van een
%% requirements-analyse, opstellen long-list (bij vergelijkende studie),
%% selectie van geschikte tools (bij vergelijkende studie, "short-list"),
%% opzetten testopstelling/PoC, uitvoeren testen en verzamelen
%% van resultaten, analyse van resultaten, ...
%%
%% !!!!! LET OP !!!!!
%%
%% Het is uitdrukkelijk NIET de bedoeling dat je het grootste deel van de corpus
%% van je bachelorproef in dit hoofstuk verwerkt! Dit hoofdstuk is eerder een
%% kort overzicht van je plan van aanpak.
%%
%% Maak voor elke fase (behalve het literatuuronderzoek) een NIEUW HOOFDSTUK aan
%% en geef het een gepaste titel.

